\section{GPU Radar Processing Performance}
\todo{Maybe write about other papers that broadly explore GPU acceleration in the context of any sensors or in the UAV context}
\subsection{GPU-accelerated radar signal processing}
\todo{Develop this futher}
Studies that are less UAV focused, such as the GPU-accelerated passive bistatic radar pipeline of Zhao et al. \cite{rs15225421} and the real-time GPU radar processing thesis of Zaknoun \cite{zaknoun2024_realtime_gpu_accelerated_radar_processing} provide more comprehensive analyses of radar software architecture and GPU implementation considerations. Zhao et al. implement compute unified device architecture (CUDA) acceleration for key stages, reporting up to a 113.13$\times$ speedup over a CPU baseline. The research is limited to only three algorithms (ECA-B clutter suppression, range–Doppler processing, and 2D CA-CFAR detection), leaving gaps for similar research that applies to other common radar algorithms such as angle sensing, coordinate transforms, kalman filters, and more. Additionally, the research did not explicitly discuss the time cost of memory reallocation and thread synchronisation. The proposed project will fill these gaps from this study by benchmarking a more comprehensive selection of tracking algorithms for CUDA implementations, and profiling bottlenecks such as host–device data transfer and branching delay to provide a clearer understanding of the benefits of parallelised processing. 
